\section*{الفصل \ref{ch:g12:photosynthesis} --- التركيب الضوئي}

\begin{solution}{exo:g12:photosynthesis:1}
\omterm{def:g10:cells-common-unit:organelle}{عضية} ذات غشاء مضاعف فيها حشوة سائلة تحوي أكداسًا من أغشية \omterm{def:g12:photosynthesis:chloroplast}{التيلاكويد}.
والأصبغة \omterm{def:g10:chemistry-of-life:families}{والبروتينات} المحوِّلة للضوء في أغشية \omterm{def:g12:photosynthesis:chloroplast}{التيلاكويد}؛ أما \omterm{def:g11:enzymes-and-phenotype:enzyme}{الإنزيمات}
التي تبني السكر ففي الحشوة.
\end{solution}

\begin{solution}{exo:g12:photosynthesis:2}
اليخضور يمتص الأزرق والأحمر ويعكس الأخضر، وهو ما نراه. والضوءان
الأزرق والأحمر يقودان \omterm{def:g10:cell-metabolism:autotrophy}{التركيب الضوئي} خير قيادة.
\end{solution}

\begin{solution}{exo:g12:photosynthesis:3}
ATP وNADP المرجَع والأكسجين؛ والأكسجين هو الناتج الثانوي، ويتحرر.
\end{solution}

\begin{solution}{exo:g12:photosynthesis:4}
من الماء المشقوق لا من ثنائي أكسيد الكربون: فالطحالب التي أعطيت ماء
موسومًا بأكسجين ثقيل حررت $\mathrm{O_2}$ ثقيلًا؛ وحين أعطيت ثنائي
أكسيد كربون موسومًا لم تفعل.
\end{solution}

\begin{solution}{exo:g12:photosynthesis:5}
ستة جزيئات $\mathrm{CO_2}$، و18 ATP، و12 NADP مرجَعًا.
\end{solution}

\begin{solution}{exo:g12:photosynthesis:6}
عند \qty{550}{nm}: الامتصاص نحو 8\%، والمعدل نحو 25\%. وعند
\qty{665}{nm}: الامتصاص نحو 90\%، والمعدل نحو 95\%. والمصابيح الخضر
تسلّم الضوء الذي تستعمله الورقة أقل ما تستعمل؛ فيوفر البيت الزجاجي
الطاقة بألا ينبت شيئًا تقريبًا.
\end{solution}

\begin{solution}{exo:g12:photosynthesis:7}
الجراثيم تسبح نحو الأكسجين؛ فتتجمع حيث يتحرر الأكسجين، أي حيث يقوم
الطحلب \omterm{def:g10:cell-metabolism:autotrophy}{بالتركيب الضوئي}؛ وذلك تحت الجزأين الأزرق والأحمر من الطيف.
\omterm{def:g10:cell-metabolism:autotrophy}{فالتركيب الضوئي}، مقيسًا بأكسجينه، يتوقف على لون الضوء ويتبع امتصاص
الأصبغة.
\end{solution}

\begin{solution}{exo:g12:photosynthesis:8}
يتحرر الأكسجين من غير أن يثبَّت أي كربون: فالمرحلة الضوئية، وهي تشق
الماء وتنقل الإلكترونات، منفصلة عن مرحلة الكربون ولا تحتاج
$\mathrm{CO_2}$؛ وإنما تحتاج ضوءًا وماء وشيئًا يتلقى الإلكترونات.
\end{solution}

\begin{solution}{exo:g12:photosynthesis:9}
من غير ضوء لا ATP ولا NADP مرجَع، فلا يستطيع الحمض الثلاثي الكربون أن
يرجَع فيتراكم؛ والمتلقي يظل يلتقط $\mathrm{CO_2}$ مدة لكنه لم يعد
يجدَّد، فيستنفد.
\end{solution}

\begin{solution}{exo:g12:photosynthesis:10}
شدة الضوء تهبط مع البعد (كمربعه تقريبًا: أي أربعة أضعاف أقل عند ضعف
البعد)، فتمشي المرحلة الضوئية أبطأ. وفي الظلام لا يتحرر أكسجين ---
بل يستهلك النبات بعضه بالتنفس.
\end{solution}

\begin{solution}{exo:g12:photosynthesis:11}
$8/44 \approx \qty{0.18}{mol}$ من $\mathrm{CO_2}$: أي $0.18/6 \approx
\qty{0.030}{mol}$ من الغلوكوز، نحو \qty{5.5}{g}؛ و\qty{0.18}{mol} من
$\mathrm{O_2}$، نحو \qty{5.8}{g}.
\end{solution}

\begin{solution}{exo:g12:photosynthesis:12}
هي لا تستعمل الضوء بنفسها، بل نواتج المرحلة الضوئية لا غير. لكنها في
نبات يبقى في الظلام تتوقف في غضون دقيقة، ما إن ينفد ATP وNADP المرجَع:
فهي تمشي في الضوء، تغذيها المرحلة الضوئية، و"المظلمة" تصف ما تحتاجه لا
متى تقع.
\end{solution}

\begin{solution}{exo:g12:photosynthesis:13}
$^{18}$O: في $\mathrm{O_2}$ المتحرر في غضون دقيقة، وهو هناك بعد ساعة
وبعد يوم (وبعضه أيضًا في ماء يصنعه التنفس).
و$^{14}$C: بعد دقيقة في الحمض الثلاثي الكربون وفي سكاكر الدورة؛ وبعد
ساعة في السكروز والنشاء؛ وبعد يوم في السيللوز والدسم \omterm{def:g10:chemistry-of-life:families}{والحموض الأمينية}،
وفي $\mathrm{CO_2}$ الذي يزفر.
\end{solution}

\begin{solution}{exo:g12:photosynthesis:14}
$20/1000 = 2\%$. والباقي ينعكس (الأخضر)، أو يمتص في أطوال موجية غير
مناسبة فيتحول حرارة، أو ينفق في النتح الذي يبخر الماء، أو يفقد حرارةً
في التفاعلات الضوئية نفسها.
\end{solution}

\begin{solution}{exo:g12:photosynthesis:15}
التنفس والاحتراق وأكسدة الصخور والحديد تستهلك الأكسجين استهلاكًا
متصلًا؛ ومن غير \omterm{def:g10:cell-metabolism:autotrophy}{تركيب ضوئي} يجدده، كان مخزون الجو سيستنفد في بضعة آلاف
من السنين. ولم يكن في الجو الباكر أكسجين حر؛ ولم يتراكم إلا بعد أن
أنتجته الكائنات التركيبية الضوئية زمنًا طويلًا جدًا --- فأول ملياري سنة
كانت فقيرة بالأكسجين.
\end{solution}

\begin{solution}{pb:g12:photosynthesis:1}
\textbf{1.} $8/44 \approx \qty{0.18}{mol}$.

\textbf{2.} $0.18/6 = \qty{0.030}{mol}$، أي نحو \qty{5.5}{g}.

\textbf{3.} \qty{0.18}{mol} من $\mathrm{O_2}$، أي نحو \qty{4.4}{L}.

\textbf{4.} جزيئا ماء لكل $\mathrm{O_2}$: أي \qty{0.36}{mol}،
ونحو \qty{6.5}{g} --- أي واحد في المئة من \qty{600}{mL} المنتوحة.

\textbf{5.} من جزيئات الماء المشقوقة في \omterm{def:g12:photosynthesis:chloroplast}{التيلاكويدات}: وقد بُيّن ذلك
بوسم الأكسجين الثقيل، وهو يظهر في $\mathrm{O_2}$ حين يوسم الماء لا
غير.

\textbf{6.} دورة واحدة لكل $\mathrm{CO_2}$: أي \qty{0.18}{mol} من الدورات.

\textbf{7.} ثلاثة ATP واثنان NADP مرجَع في الدورة: أي \qty{0.55}{mol}
من ATP و\qty{0.36}{mol} من NADP المرجَع.

\textbf{8.} كل ماء يشق يعطي إلكترونين، أي NADP مرجَعًا واحدًا:
\qty{0.36}{mol} من الماء --- وهو الرقم نفسه في السؤال 4، كما يجب أن
يكون.

\textbf{9.} يتوقف الأكسجين في الحال؛ ويكف NADP المرجَع وATP عن التكون
وينفدان في غضون ثوان؛ ويتوقف السكر ما إن ينفد الحاملان.

\textbf{10.} لم تعد \omterm{prop:g12:photosynthesis:calvin}{دورة كالفن} تنفق الحاملين؛ فمن غير شيء يتلقى
إلكتروناتهما وطاقتهما، تتكدس المرحلة الضوئية وتبطئ في غضون ثوان ---
فالمرحلتان مقترنتان بالحاملين.

\textbf{11.} $0.030 \times 2870 \approx \qty{86}{kJ}$ في الساعة، أي
$86000/3600 \approx \qty{24}{W}$.

\textbf{12.} $24/1000 = 2.4\%$ من ضوء الشمس الكامل؛ و$24/450 \approx
5\%$ من الأطوال الموجية الممتصة.

\textbf{13.} $8 \times 176 \approx \qty{1400}{kJ}$ من الفوتونات لكل
مول من $\mathrm{CO_2}$؛ والمخزون: $2870/6 \approx \qty{480}{kJ}$.

\textbf{14.} نحو الثلث؛ والباقي يتحرر حرارةً في نقلات الطاقة المتعاقبة
في المرحلتين.

\textbf{15.} المكسب الصافي $5.5 - 0.5 = \qty{5}{g}$ من الغلوكوز في
الساعة. ويقاس بأخذ $\mathrm{CO_2}$ الصافي في الضوء؛ أما التنفس وحده
فيقاس في الظلام ويضاف لنحصل على \omterm{def:g10:cell-metabolism:autotrophy}{التركيب الضوئي} الإجمالي.

\textbf{16.} $\num{8.5e11}/\num{1e11} \approx 8.5$ سنوات --- ومع
تثبيت طحالب المحيط قدر ذلك، نحو 4 سنوات.

\textbf{17.} التنفس (عند النبات والحيوان والفطور والجراثيم)، مع
الحريق، يرد ثنائي أكسيد الكربون بالمعدل نفسه تقريبًا؛ فتوازن التثبيت
والرد، وظل المخزون ثابتًا حتى أضاف حرق الوقود الأحفوري مصدرًا جديدًا.

\textbf{18.} \qty{1e11}{t} من الكربون هي $\num{8.3e12}\,\unit{mol}$؛
والعدد نفسه من مولات $\mathrm{O_2}$ نحو \qty{2.7e11}{t} من
الأكسجين.

\textbf{19.} $\num{1.2e15}/\num{2.7e11} \approx 4500$ سنة من الإنتاج
(ونحو 2000 مع حصة المحيط). لكن التنفس يستهلك الأكسجين بمعدل يقارب
معدل صنعه \omterm{def:g10:cell-metabolism:autotrophy}{بالتركيب الضوئي}، فالمخزون توازن لا تراكم: والأكسجين الذي في
الهواء بني على مدى مئات الملايين من السنين بالفائض الصغير للإنتاج على
الاستهلاك.

\textbf{20.} نحو 2\% من ضوء الشمس الكامل (أي 5\% من الأطوال الموجية
الصالحة) يخزن سكرًا؛ والأكسجين المتحرر هو أكسجين الماء المشقوق؛ ونحو
ربع كربون الجو يمر عبر \omterm{def:g10:cell-metabolism:autotrophy}{التركيب الضوئي} كل سنة.
\end{solution}
